Die Erstellung einer Bibliografie erfolgt in mehreren Schritten.
Listing~\ref{lst:latex:bibliography:template} zeigt den grundlegenden Aufbau.

\begin{compactenum}
\item Zunächst muss in der Präambel das Paket~\ctanpkg{biblatex} geladen
      werden.
\item Anschließend werden noch im Vorspann mit
      \lstinline[language={[LaTeX]{TeX}}]+\addbibresource+
      die Literaturdatenbanken geladen. Die Vorlage verwendet hierfür
      standardmäßig eine einzige Datei~\texttt{literatur.bib}. Für umfangreiche
      Arbeiten oder komplexe Literaturstudien ist es auch möglich, die
      Bibliografie aus mehreren \texttt{.bib}-Dateien zusammenzusetzen.
      Zweckmäßiger Weise legt man diese Dateien dann in einem separaten
      Unterverzeichnis ab. Jede Datei mit einem entsprechenden
      \lstinline[language={[LaTeX]{TeX}}]+\addbibresource+-Befehl
      geladen werden.
\item Die eigentliche Bibliografie wird innerhalb des Dokuments mit
      \lstinline[language={[LaTeX]{TeX}}]+\printbibliography+
      ausgegeben.
\end{compactenum}

\lstinputlisting[language={[LaTeX]{TeX}},style=block,caption={Grundgerüst einer Bibliografie},label={lst:latex:bibliography:template}]{code/latex_bibliography_template.de.tex}

Standardmäßig erscheint die Bibliografie nicht im Inhaltsverzeichnis.
Listing~\ref{lst:latex:bibliography:toc} enthält die notwendigen Optionen, mit
denen sich die Bibliografie im Inhaltsverzeichnis verankern lässt.

\lstinputlisting[language={[LaTeX]{TeX}},style=block,caption={Bibliografie im Inhaltsverzeichnis aufführen},label={lst:latex:bibliography:toc}]{code/latex_bibliography_toc.de.tex}

Sowohl das Paket als auch die oben genannten Befehle lassen sich durch eine
Vielzahl von weiteren Optionen angepassen. Hierfür sei auf die entsprechende
Dokumentation von \cite{pkg_biblatex} verwiesen.

Der Quellenverweis erfolgt mit dem
\lstinline[language={[LaTeX]{TeX}}]+\cite+-Befehl. Das Schlüsselwort für die
Fundstelle entspricht dem Schlüsselwort, welches in der \hologo{BibTeX}-Datei
vergeben wurde. Daneben gibt es noch eine Vielzahl weiterer Zitierbefehle, die
für verschiedene Anweungsfälle gedacht sind. Der folgende Absatz listet einige
Varianten beispielhaft auf. Der \LaTeX-Code hierzu ist in
Listing~\ref{lst:latex:bibliography:cite} abgedruckt. Eine genaue Beschreibung
jeder Variante würde den Rahmen dieser Anleitung sprengen. Hierfür sei auf
\cite{pkg_biblatex,rees_biblatex_cheat} verwiesen.

\input{code/latex_bibliography_cite.de.tex}

\lstinputlisting[language={[LaTeX]{TeX}},style=block,caption={Literaturverweise},label={lst:latex:bibliography:cite},morekeywords={footcite,citeauthor,citetitle,fullcite,footfullcite}]{code/latex_bibliography_cite.de.tex}
